\section{Verification Framework for Subnormal Possibility Distributions}\label{app:verification_framework}

This appendix formalizes the verification metrics for \clyfar{}'s native possibility outputs (Product~3). All definitions use the notation of Section~\ref{sec:poss_framework}: $\pi$ denotes the raw (subnormal) possibility distribution, $m = \max_{\omega \in \Omega} \pi(\omega)$ is the subnormality, and $\ign = 1 - m$ is ignorance (Eq.~\eqref{eq:ignorance}). The tripartite pignistic bridge and information-gain verification (Appendix~\ref{app:information_theory}) are not repeated here; they evaluate Product~2 (probability-converted forecasts) in a separate verification lane.

\subsection{Normalization Protocol}\label{sec:norm_protocol}

\clyfar{} outputs subnormal possibility distributions; some verification quantities require the raw distribution $\pi$ while others require the normalized form $\pinorm = \pi / m$. Two normalization operations serve distinct purposes:

\begin{enumerate}[leftmargin=2em]
\item \textbf{Shape normalization} ($\pi \to \pinorm$): Rescales so that $\max(\pinorm) = 1$, isolating the model's relative preferences among outcomes. Used for computing \Nc{} (Eq.~\eqref{eq:cond_nec}) and the native scorecard (Section~\ref{sec:native_scorecard}) within the possibility verification lane.
\item \textbf{Tripartite pignistic bridge} ($\pi \to$ probability vector): Converts possibilities to probabilities with an explicit ignorance outcome (Section~\ref{sec:poss_prob_bridge}). Used for information-gain verification in the probability lane.
\end{enumerate}

These are not alternatives: normalization operates within the possibility space; the pignistic bridge converts to the probability space.

Table~\ref{tab:raw_norm} specifies which form each quantity uses.

\begin{table}[htbp]
\centering
\caption{Raw vs.\ normalized possibility assignments for verification quantities.}
\label{tab:raw_norm}
\small
\begin{tabularx}{\columnwidth}{>{\raggedright\arraybackslash}p{0.20\columnwidth}ccX}
\toprule
Quantity & Symbol & Uses & Reason \\
\midrule
Ignorance & $\ign$ & $\pi$ & Gap to normality \\
Event possibility & $\poss(A_T)$ & $\pi$ & Upper plausibility bound \\
Cumulative upper bound & $U_k$ & $\pi$ & For IRLS \\
Depth-of-truth & $\alpha^*$ & $\pinorm$ & Shape quality \\
Nonspecificity & $\eta$ & $\pinorm$ & Shape diffuseness \\
Conditional necessity & $\Nc(A)$ & $\pinorm$ & Certainty given coverage \\
Lower threshold bound & $L_t$ & Both & Coverage $\times$ certainty \\
Cumulative lower bound & $L_k$ & Both & IRLS counterpart of $L_t$ \\
\bottomrule
\end{tabularx}
\end{table}

\subsection{Native Possibility Scorecard}\label{sec:native_scorecard}

Five summary statistics evaluate each forecast--observation pair, where $c_{\mathrm{obs}}$ denotes the observed category and $K = |\Omega|$ is the number of categories.

\textbf{Depth-of-truth} (compatibility of observation with model):
\begin{equation}
\alpha^* = \pinorm(c_{\mathrm{obs}}) \label{eq:alpha_star}
\end{equation}
where $\alpha^* = 1$ when the truth coincides with the model's peak category and $\alpha^* = 0$ when the model assigns zero possibility to the observed outcome.

\textbf{Nonspecificity} \citep{Klir1995-va}, measuring diffuseness of the normalized shape:
\begin{equation}
\eta = \frac{1}{K} \sum_{c=1}^{K} \pinorm(c) \label{eq:ns}
\end{equation}
ranging from $1/K$ (sharp) to $1$ (uniform).

\textbf{Resolution gap} (net resolution):
\begin{equation}
\delta = \alpha^* - \eta \label{eq:rg}
\end{equation}
Positive $\delta$ indicates the model assigned more possibility to the truth than to an average category.

\textbf{Ignorance} is $\ign = 1 - m$ (Eq.~\eqref{eq:ignorance}).

\textbf{Conditional necessity of truth}:
\begin{equation}
\Nc^* = \Nc(c_{\mathrm{obs}}) = 1 - \max_{\omega \neq c_{\mathrm{obs}}} \pinorm(\omega) \label{eq:nc_star}
\end{equation}
where $\Nc^* > 0$ only when the observed category dominates all alternatives in the normalized distribution. This is a direct application of Eq.~\eqref{eq:cond_nec} evaluated at the observation.

These five statistics map onto the three-component framework of Section~\ref{sec:three_component}: $\alpha^*$ and $\delta$ evaluate possibility-derived shape quality, $\ign$ evaluates the ignorance signal, and $\Nc^*$ evaluates conditional necessity. Aggregates $\overline{\alpha^*}$, $\overline{\eta}$, $\overline{\delta}$, $\overline{\ign}$, $\overline{\Nc^*}$ are reported with bootstrap confidence intervals.

\subsection{Interval Log Score for Threshold Events}\label{sec:interval_log}

For a threshold event $A_T \subseteq \Omega$ (e.g., $A_T = \{\text{elevated}, \text{extreme}\}$ for ozone exceedance), the possibilistic forecast produces an interval $[L, U]$ rather than a point probability. The bounds derive from the three-component framework:

\begin{align}
U &= \poss(A_T) = \max_{\omega \in A_T} \pi(\omega) \label{eq:U_bound} \\
L &= m \cdot \Nc(A_T) = m \cdot \bigl[1 - \max_{\omega \notin A_T} \pinorm(\omega)\bigr] \label{eq:L_bound}
\end{align}

The \textit{interval log score} is:
\begin{equation}
\mathrm{ILS} = e \cdot \bigl[-\log_2\bigl(\max(U,\, \epsilon)\bigr)\bigr] + (1 - e) \cdot \bigl[-\log_2\bigl(\max(1 - L,\, \epsilon)\bigr)\bigr] \label{eq:interval_log}
\end{equation}
where $e \in \{0, 1\}$ is the event indicator and $\epsilon$ is a small positive floor (herein $\epsilon = 0.01$, capping the maximum penalty at $-\log_2(0.01) \approx 6.64$~bits). When $L = U = p$ (a point probability), ILS reduces to the standard log score.

\textbf{Motivation:} Quadratic scores (Brier) treat a forecast of 0\% and 5\% almost identically for a rare event that occurs. Logarithmic scoring penalizes confident wrong forecasts far more severely, which is appropriate for rare exceedance events (${\sim}5\%$ base rate across ${\sim}960$ station-days per season). The interval log score is the natural information-theoretic extension of the interval Brier score; the latter is recoverable as a second-order Taylor approximation of ILS around $p = 0.5$ \citep{Roulston2002-eq}.

\textbf{Three-component interpretation:} The forecast assigns possibility $\poss = U$ to exceedance, conditional necessity $\Nc = L/m$, and ignorance $\ign = 1 - m$. The interval $[L, U] = [m \cdot \Nc, \poss]$ brackets the event's plausibility; width $U - L$ reflects ambiguity between ``merely possible'' and ``essentially certain.''

\textbf{Properties:}
\begin{itemize}[leftmargin=2em]
\item $0 \leq L \leq U \leq 1$ always holds.
\item Under normality ($m = 1$): $L = \nec(A_T)$, $U = \poss(A_T)$.
\item Under extreme subnormality ($m \to 0$): $U \to 0$, so ILS for an observed event approaches $-\log_2(\epsilon)$ (maximum penalty); $L \to 0$, so ILS for a non-event approaches $-\log_2(1) = 0$ (no penalty). This correctly distinguishes ``no information and the event happened'' (costly) from ``no information and nothing happened'' (costless).
\end{itemize}

\textbf{Worked example} (Example~1 from Section~\ref{sec:subnormal}, $\pi = \{\text{bkg}: 0,\;\text{mod}: 0.2,\;\text{elev}: 0.25,\;\text{ext}: 0.75\}$, $A_T = \{\text{elev}, \text{ext}\}$):\; $m = 0.75$;\; $U = 0.75$;\; $\pinorm = \{0,\; 0.267,\; 0.333,\; 1.0\}$;\; $\Nc(A_T) = 1 - 0.267 = 0.733$;\; $L = 0.75 \times 0.733 = 0.55$;\; $L \leq U$ confirmed ($0.55 \leq 0.75$).\; If exceedance occurs: $\mathrm{ILS} = -\log_2(0.75) = 0.42$~bits.\; If not: $\mathrm{ILS} = -\log_2(1 - 0.55) = -\log_2(0.45) = 1.15$~bits.

\subsection{Imprecise Ranked Logarithmic Score}\label{sec:irls}

For $K$ ordered categories, the cumulative possibility and necessity bounds are:
\begin{align}
U_k &= \max_{i \leq k}\, \pi(i) \\
L_k &= m \cdot \bigl[1 - \max_{i > k}\, \pinorm(i)\bigr]
\end{align}
with $\max$ over an empty set defined as $0$, so that $U_K = m$ and $L_K = m$.

The imprecise ranked logarithmic score is:
\begin{equation}
\mathrm{IRLS} = \frac{1}{K} \sum_{k=1}^{K} \bigl[e_k \cdot S(U_k) + (1 - e_k) \cdot S(1 - L_k)\bigr] \label{eq:irls}
\end{equation}
where $S(p) = -\log_2(\max(p, \epsilon))$ and $e_k = \mathbf{1}\{c_{\mathrm{obs}} \leq k\}$. When $L_k = U_k$ for all $k$, IRLS reduces to the standard ranked logarithmic score.

\textbf{Properties:} $U_k$ is non-decreasing in $k$ (maximum over a growing set); $L_k$ is non-decreasing in $k$ (complement of a maximum over a shrinking set); $L_k \leq U_k$ for all $k$.

\textbf{Note:} IRLS is formally defined here for completeness but its application is deferred to a follow-on verification study with a larger reforecast dataset.

\subsection{Possibilistic Climatology}\label{sec:poss_climatology}

A possibilistic baseline analogous to probabilistic climatology is needed for skill scores. The recommended construction derives consonant possibility distributions from the climatological CDF.

For continuous variables with climatological CDF $F_{\mathrm{clim}}(z)$:
\begin{equation}
\pi_{\mathrm{clim}}(z) = 1 - 2\,\bigl|F_{\mathrm{clim}}(z) - 0.5\bigr| \label{eq:poss_clim}
\end{equation}
This peaks at the climatological median ($\pi_{\mathrm{clim}} = 1$ where $F_{\mathrm{clim}} = 0.5$) and produces central nested $\alpha$-cut sets $S_{\mathrm{clim}}(\alpha) = [Q_{\mathrm{clim}}(\alpha/2),\; Q_{\mathrm{clim}}(1 - \alpha/2)]$.

For $K$ ordered categories with smoothed climatological frequencies $p_k$, the mid-rank method is recommended: compute cumulative frequency $F_{k-1} = \sum_{j < k} p_j$, category mid-rank $u_k = F_{k-1} + 0.5\, p_k$, and define $\pi_{\mathrm{clim}}(k) = 1 - 2\,|u_k - 0.5|$.

Climatological ignorance is zero by convention ($I_{\mathrm{clim}} = 0$): climatology is a reference distribution, not a live model state.

\subsection{Tripartite Value Demonstration}\label{sec:tripartite_demo}

The five-number scorecard (Section~\ref{sec:native_scorecard}) can be used to demonstrate that each of the three named components---possibility, ignorance, and conditional necessity---independently contributes verification value:

\begin{enumerate}[leftmargin=2em]
\item \textbf{Possibility compatibility} ($\alpha^*$, $\delta$): Does positive resolution gap persist when comparing \clyfar{} against persistence or climatology baselines?
\item \textbf{Ignorance self-awareness} ($\ign$): Is $\overline{\delta}$ lower in high-ignorance cases than in low-ignorance cases? A positive answer confirms that the model's self-reported uncertainty tracks actual forecast difficulty.
\item \textbf{Conditional necessity reliability} ($\Nc^*$): When $\Nc(c) > \tau$, does the conditional hit rate exceed the base rate of category $c$? This demonstrates that \Nc{} identifies cases where the model is not merely plausible but confidently correct.
\end{enumerate}

At pilot sample sizes ($\sim$90 station-days), these tests are supplemented by 3--4 illustrative case studies showing how the full $(\poss, \ign, \Nc)$ triplet conveys information invisible to the defuzzified scalar alone.
